\PassOptionsToPackage{dvipsnames,svgnames,table}{xcolor}
\documentclass[11pt]{article}
\usepackage[a4paper,margin=22mm]{geometry}
\usepackage[no-math]{fontspec}
\setmainfont{DejaVu Serif}
\usepackage{amsmath,amssymb,amsthm,mathtools,graphicx,xcolor,hyperref,xurl,xspace,ifthen,enumerate}
\usepackage[shortlabels]{enumitem}
\setlength{\emergencystretch}{4em}
\AtBeginDocument{\special{pdf:mapline pzdr Dingbats <uzdr.pfb}}
\SetMathAlphabet{\mathbf}{normal}{OT1}{cmr}{bx}{n}
\SetMathAlphabet{\mathbf}{bold}{OT1}{cmr}{bx}{n}
\newtheorem{theo}{Theorem}[section]
\newtheorem{lemm}[theo]{Lemma}
\newtheorem{thm}[theo]{Theorem}
\newtheorem{proposition}[theo]{Proposition}
\newtheorem{corollary}[theo]{Corollary}
\newtheorem{definition}[theo]{Definition}
\newtheorem{theorem}[theo]{Theorem}
\newtheorem*{remas*}{Remarks}
\newtheorem{remas}[theo]{Remarks}
\newtheorem*{exam*}{Example}
\newtheorem{ques}[theo]{Question}
\providecommand{\og}{«\nobreakspace}
\providecommand{\fg}{\nobreakspace»}
\newtheorem{lem}[theo]{Lemma}
\newtheorem{lemma}[theo]{Lemma}
\newtheorem{prop}[theo]{Proposition}
\newtheorem{coro}[theo]{Corollary}
\newtheorem{corol}[theo]{Corollary}
\newtheorem*{conj*}{Conjecture}
\newtheorem{conj}[theo]{Conjecture}
\newtheorem{conjecture}[theo]{Conjecture}
\newtheorem{Proposition}[theo]{Proposition}
\newtheorem{theoreme}[theo]{Theorem}
\theoremstyle{definition}
\newtheorem{defi}[theo]{Definition}
\newtheorem{exam}[theo]{Example}
\newtheorem{exem}[theo]{Example}
\newtheorem{remark}[theo]{Remark}
\newtheorem{example}[theo]{Example}
\newtheorem{exams}[theo]{Examples}
\newtheorem{rema}[theo]{Remark}
\newtheorem{nota}[theo]{Notation}
\newtheorem{notation}[theo]{Notation}
\newtheorem*{notas}{Notations}
\newtheorem{result}[theo]{Result}
\newtheorem{question}[theo]{Question}
\newtheorem{prob}[theo]{Problem}
\newtheorem*{rema*}{Remark}
\newtheorem*{defi*}{Definition}
\newtheorem*{theo*}{Theorem}
\newtheorem*{lemm*}{Lemma}
\newtheorem*{prop*}{Proposition}
\newcommand{\enoncename}{}
\newtheorem{corpusenonce}[theo]{\enoncename}
\newenvironment{enonce}[2][]{\renewcommand{\enoncename}{#2}\begin{corpusenonce}}{\end{corpusenonce}}
\newenvironment{enonce*}[2][]{\par\medskip\noindent\textbf{#2.}\itshape}{\par\medskip}
\newenvironment{proofc}[1][\proofname]{\begin{proof}[#1]}{\end{proof}}
\newcommand{\Mc}{\mathcal{M}}
\newcommand{\pointir}{.\ ---\ }
\newcommand{\equalenv}[2]{\expandafter\let\csname#1\expandafter\endcsname\csname#2\endcsname\expandafter\let\csname end#1\expandafter\endcsname\csname end#2\endcsname}
\providecommand{\diagbox}[3][]{\begin{tikzpicture}[x=1em,y=1em,baseline=(current bounding box.center)]\useasboundingbox(0,0) rectangle(3,1.4);\draw(0,1.4)--(3,0);\node at(.5,.3){#2};\node at(2.5,1.1){#3};\end{tikzpicture}}
\providecommand{\eg}{e.g.\ }
\providecommand{\ie}{i.e.\ }
\providecommand{\cf}{cf.\ }
\providecommand{\longthanks}[1]{\par\smallskip\noindent #1\par\smallskip}
\providecommand{\qbinom}[2]{\genfrac{[}{]}{0pt}{}{#1}{#2}}

\providecommand{\mathof}[1]{\mathrm{#1}}

\numberwithin{equation}{section}


\def\endappendix{}


\usepackage{latexsym,amsmath,amssymb,verbatim,tikz,tikz-cd}
\usepackage{mathtools}
\usepackage{graphicx,xcolor}

\usepackage[english]{babel}
\usepackage{amssymb}
\usepackage{amsmath}
\usepackage{amsthm}
\usepackage{graphicx}
\usepackage{csquotes}
\usepackage{bm}
\usepackage{enumitem}
\usepackage{subcaption}
\usepackage{caption}
\renewcommand{\thesubfigure}{\Alph{subfigure}}
\usepackage{dirtytalk}
\providecommand{\nopunct}{\spacefactor1007\relax}
\newtheorem{problem}[theo]{Problem}

\newcommand{\sym}{\mathrm{Sym}}
\newcommand{\qsym}{\mathrm{QSym}}
\newcommand{\Des}{\operatorname{Des}}
\newcommand{\Asc}{\operatorname{Asc}}
\newcommand{\Syt}{\operatorname{SYT}}
\newcommand{\Short}{\operatorname{Short}}
\newcommand{\Loop}{\operatorname{Loop}}
\newcommand{\Stable}{\operatorname{Stable}}
\newcommand{\len}{\operatorname{len}}
\newcommand{\height}{\operatorname{height}}
\newcommand{\core}{\operatorname{core}}
\newcommand{\res}{\operatorname{res}}
\newcommand{\row}{\operatorname{row}}
\newcommand{\insrt}{\operatorname{insert}}
\newcommand{\UP}{\operatorname{UP}}
\newcommand{\DOWN}{\operatorname{DOWN}}


\newcommand{\ZZ}{\mathbb{Z}}
\newcommand{\NN}{\mathbb{N}}
\newcommand{\RR}{\mathbb{R}}
\newcommand{\QQ}{\mathbb{Q}}
\newcommand{\PP}{\mathbb{P}}

\newcommand{\F}{\mathcal{F}}
\newcommand{\Q}{\mathcal{Q}}
\newcommand{\A}{\mathcal{A}}
\newcommand{\M}{\mathcal{M}}
\newcommand{\I}{\mathcal{I}}
\newcommand{\PM}{\mathcal{PM}}
\newcommand{\Pcal}{\mathcal{P}}


\newcommand{\odds}{\mathrm{odds}}

\makeatletter
\newcommand{\proofstep}[1]{
  \par
  \addvspace{\medskipamount}
  \textbf{#1\@addpunct{.}}\enspace\ignorespaces
}

\newcommand\restr[2]{{
  \left.\kern-\nulldelimiterspace 
  #1 
  \vphantom{\big|} 
  \right|_{#2} 
  }}
\makeatother












\usepackage{cleveref}
\crefname{originaltheorem}{original theorem}{original theorems}
\Crefname{originaltheorem}{Original theorem}{Original theorems}
\makeatletter
\newlabel{thm:criterion schur positive young}{{2.5}{source}{}{sourcepointer-1}{}}
\newlabel{thm:criterion schur positive young@cref}{{[originaltheorem][0][]2.5}{[1][0][]source}{}{}{}}
\makeatother
\title{Hook-shaped Schur expansions are equivalent to uniformity of the descent-set distribution}
\author{Avichai Marmor}
\date{doi:10.5802/alco.351; CC BY 4.0}
\begin{document}
\maketitle
\noindent\textit{Editorial note.} The counted statement is the author's labelled proposition `prop:hook schur positive` (source lines 1334-1341); the proof environment begins on the next line after \textbackslash{}end\{prop\} apart from one blank line, so the pairing is by adjacency.\par\medskip
\noindent\textit{Editorial note.} The retained context is the author's own lead-in sentence introducing the criterion.\par\medskip
\noindent\textit{Editorial note.} The retained preamble is the editorially prepared standalone preamble retained for this source, and the retained bibliography of the source is bundled because the retained proof cites a reference for the converse direction.\par\medskip
\noindent\textit{Editorial note.} One external reference records the Schur-positivity criterion invoked by label inside the retained proof; it is a navigation pointer only.\par\medskip
\noindent\textit{Editorial scope.} One original result of the pinned publisher article, counted once: the hook-shaped criterion for Schur-positivity (author label prop:hook schur positive).\par\medskip
\noindent\textit{Allowed standard prerequisites (editorial).}\begin{itemize}
\item Hook Schur function. s\_\{N-k,1\textasciicircum{}k\} is the Schur function of the hook shape with k boxes below the first row.. It is the shape appearing in the claimed expansion.
\item Descent set of a hook tableau. Every standard Young tableau of hook shape (N-k,1\textasciicircum{}k) has exactly k descents, and each k-subset of [N-1] arises from exactly one such tableau.. It is the bijective step in the forward direction at source line 1344.
\item Schur-positivity criterion. A set is symmetric for D if and only if its generating function expands over standard Young tableaux with coefficients depending only on the shape.. It converts the displayed generating-function identity into the Schur expansion at source line 1349.
\end{itemize}\medskip
\ifdefined\hypertarget\hypertarget{sourcepointer-1}{}\fi
\noindent\textit{Source reference.} Original theorem 2.5: sources/alco-351-matching-schur-2024/publisher-source.tex, source lines 368--368. Original criterion of Adin and Roichman invoked at source line 1349. It is a navigation pointer: the counted proposition's own proof is retained line for line and nothing is replaced by this reference. The reference number was inferred by counting the theorem-class environments of that section of the pinned source; it is not confirmed by the retained bibliography or by a label-number dump of the original, so it must be read as a navigation pointer only.\par\medskip
There is a simple well-known criterion for sets under which the generating function $\Q_{D}(\A)$ is non-negatively spanned by the Schur functions of hook shape $s_{N-k,1^k}$:

\begin{proposition}[folklore]
\label{prop:hook schur positive}
    A set $\A$ with a statistic $D:\A \to 2^{[N-1]}$ has a generating function of the form
    \[
        \Q_{D}(\A) = \sum_{k=0}^{N-1} c_k s_{N-k,1^k},\ c_k \in \NN
    \]
    if and only if for every $J \subseteq [N-1]$, $|\{a \in \A \mid D(a) = J\}|$ depends only on $|J|$. In that case, $c_k = |\{a \in \A \mid D(a) = [k]\}|$.
\end{proposition}

\begin{proof}
    On one hand, assume that there exist $c_k \in \NN$ such that $|\A(D=J)| = c_{|J|}$ for all $J \subseteq [N-1]$. Notice that every tableau of a hook shape $T \in \Syt(N-k,1^k)$ has $|\Des(T)| = k$. Conversely, for every set $J \subseteq [N-1]$ with $|J| = k$ there exists a unique tableau $T \in \Syt(N-k,1^k)$ such that $\Des(T) = J$. Therefore, we obtain the equation
    \[
        \sum_{a \in \A} \boldsymbol{t}^{D(a)} = \sum_{k=0}^{N-1} c_k \sum_{T \in \Syt(N-k,1^k)} \boldsymbol{t}^{\Des(T)},
    \]
    where $\boldsymbol{t}^J := \prod_{j \in J} t_j$ for $J \subseteq [N-1]$.
    By Theorem~\ref{thm:criterion schur positive young}, this equation implies the statement.

    The other direction of the proposition states that if $\Q_{D}(\A) = \sum_{k=0}^{N-1} c_k s_{N-k,1^k}$, then $|\{a \in \A \mid D(a) = J\}| = c_{|J|}$. Adin and Roichman \cite[Corollary 8.1]{adin_roichman_fine_set} established that a generating function uniquely determines the descent set distribution, which implies the statement.
\end{proof}

\begin{thebibliography}{MMMMMM}
\bibitem[3]{adin_roichman_fine_set} Adin, Ron M.; Roichman, Yuval Matrices, characters and descents, Linear Algebra Appl., Volume 469 (2015), pp. 381-418 \href{https://alco.centre-mersenne.org/articles/10.5802/alco.351/#adin_roichman_fine_set}{Publisher reference}.
\bibitem[13]{knuth_classes_schur_positive} Gessel, Ira M. Multipartite P-partitions and inner products of skew Schur functions, Combinatorics and algebra (Boulder, Colo., 1983) (Contemp. Math.), Volume 34, Amer. Math. Soc., Providence, RI, 1984, pp. 289-317 \href{https://alco.centre-mersenne.org/articles/10.5802/alco.351/#knuth_classes_schur_positive}{Publisher reference}.
\bibitem[14]{gessel1993counting} Gessel, Ira M.; Reutenauer, Christophe Counting permutations with given cycle structure and descent set, J. Combin. Theory Ser. A, Volume 64 (1993) no. 2, pp. 189-215 \href{https://alco.centre-mersenne.org/articles/10.5802/alco.351/#gessel1993counting}{Publisher reference}.
\bibitem[16]{S3_patterns_list} Hamaker, Zachary; Pawlowski, Brendan; Sagan, Bruce E. Pattern avoidance and quasisymmetric functions, Algebr. Comb., Volume 3 (2020) no. 2, pp. 365-388 \href{https://alco.centre-mersenne.org/articles/10.5802/alco.351/#S3_patterns_list}{Publisher reference}.
\bibitem[26]{sagan_book} Sagan, Bruce E. The symmetric group: Representations, combinatorial algorithms, and symmetric functions, Graduate Texts in Mathematics, 203, Springer-Verlag, New York, 2001, xvi+238 pages \href{https://alco.centre-mersenne.org/articles/10.5802/alco.351/#sagan_book}{Publisher reference}.
\end{thebibliography}
\end{document}
